\begin{figure}[t]
\centering
\begin{tikzpicture}[
  font=\scriptsize,
  blk/.style={rounded corners=2pt, draw=black!70, fill=blue!7, minimum width=0.95cm, minimum height=0.42cm, align=center},
  tsk/.style={rounded corners=2pt, draw=black!80, fill=black!8, minimum width=0.9cm, minimum height=0.4cm, align=center},
  call/.style={circle, draw=black!70, fill=orange!16, minimum size=0.42cm, inner sep=0.5pt, align=center},
  art/.style={rounded corners=1pt, draw=black!70, fill=green!12, minimum width=0.95cm, minimum height=0.4cm, align=center},
  st/.style={rounded corners=1pt, draw=black!70, fill=purple!12, minimum width=0.9cm, minimum height=0.4cm, align=center},
  fail/.style={circle, draw=red!70, fill=red!14, minimum size=0.42cm, inner sep=0.5pt, align=center},
  rel/.style={->, >=Stealth, draw=black!55, thin},
  faint/.style={->, >=Stealth, draw=black!30, thin},
]
% task
\node[tsk] (t) at (5.3,3.3) {task};
% blocks
\node[blk] (b1) at (0,2.1) {$P_1$};
\node[blk] (b2) at (2.65,2.1) {$P_2$};
\node[blk] (b3) at (5.3,2.1) {$P_3$};
\node[blk] (b4) at (7.95,2.1) {$P_4$};
\node[blk] (b5) at (10.6,2.1) {$P_5$};
\foreach \b in {b1,b2,b3,b4,b5}{\draw[faint] (t)--(\b);}

% B1: read -> artifact a1
\node[call] (c1) at (0,1.15) {\tiny r};
\node[art]  (a1) at (0,0.25) {\tiny $v_1$};
\draw[rel] (b1)--(c1) node[midway,right,font=\tiny]{exec};
\draw[rel] (c1)--(a1) node[midway,right,font=\tiny]{prod};

% B2: write -> state s1
\node[call] (c2) at (2.65,1.15) {\tiny w};
\node[st]   (s1) at (2.65,0.25) {\tiny state};
\draw[rel] (b2)--(c2);
\draw[rel] (c2)--(s1) node[midway,right,font=\tiny]{mut};

% B3: read -> artifact a2, equiv to a1  (Case 1)
\node[call] (c3) at (5.3,1.15) {\tiny r};
\node[art]  (a2) at (5.3,0.25) {\tiny $v_2$};
\draw[rel] (b3)--(c3);
\draw[rel] (c3)--(a2);
\draw[->, >=Stealth, dotted, green!45!black, thick] (a2) to[out=210,in=-30] node[midway,above=1pt,font=\tiny,text=green!45!black]{equiv} (a1);
\node[font=\tiny, text=green!45!black] at (2.65,-0.9) {(1) skip re-query};

% B4: failure + PATCH  (Case 2)
\node[fail] (f) at (7.95,1.15) {\tiny F};
\draw[rel,red!60!black] (b4)--(f) node[midway,right,font=\tiny,text=red!60!black]{fails\_at};
\draw[->, >=Stealth, red!65!black, thick] (f) to[out=150,in=210,looseness=6] node[midway,left,font=\tiny,text=red!65!black]{PATCH} (b4.west);
\node[font=\tiny, text=red!65!black] at (7.95,-0.9) {(2) targeted retry};

% B5: reads earlier state via supersedes  (Case 3)
\draw[->, >=Stealth, purple!60!black, thick] (s1) to[out=-70,in=-90] node[pos=0.85,above,font=\tiny,text=purple!60!black]{reads / supersedes} (b5);
\node[font=\tiny, text=purple!60!black] at (10.6,-0.9) {(3) reuse write};

\end{tikzpicture}
\caption{An example execution graph built by GraphPTC. Blocks $P_t$ are projected into typed calls (read $r$, write $w$), artifacts $v$, states, and failures $F$. (1) block $P_3$ produces $v_2$ that an \texttt{equivalent\_to} edge links to the earlier $v_1$, so the agent skips a redundant query. (2) a failed call in $P_4$ adds a \texttt{FAILURE} the agent repairs with a targeted \textsc{Patch} instead of re-running the block. (3) block $P_5$ reads the state written by $P_2$ through the \texttt{supersedes} chain instead of re-deriving it.}
\label{fig:casegraph}
\end{figure}
